\resizebox{1.0\textwidth}{!}{
\begin{tikzpicture}[node distance=0.9cm]

  % ── 节点样式定义（局部） ──────────────────────────
  \tikzset{
    mybox/.style={
      draw=blue!60, fill=blue!8, rounded corners=4pt,
      text width=6.0cm, align=center, font=\normalsize,
      minimum height=0.9cm, inner sep=6pt
    },
    mydiamond/.style={
      draw=orange!80, fill=orange!15,
      diamond, aspect=2.2,
      align=center, font=\normalsize, inner sep=3pt
    },
    sidegreen/.style={
      draw=green!60!black, fill=green!10, rounded corners=4pt,
      text width=3.8cm, align=left, font=\normalsize,
      inner sep=6pt
    },
    sideleft/.style={
      draw=blue!40, fill=blue!5, rounded corners=4pt,
      text width=2.8cm, align=center, font=\normalsize,
      inner sep=5pt
    },
    myarrow/.style={->, >=stealth, thick},
    greenarrow/.style={->, >=stealth, thick, green!60!black}
  }

  % ── 主流程节点 ───────────────────────────────────
  \node[mybox] (init) {
    \textbf{初始化：计算时间 $t = 0$}\\[2pt]
    设定初始单元平均值 $\bar{u}$\\
    （密度、速度、压力等）
  };

  \node[mydiamond, below=of init] (start) {开始计算};

  \node[mybox, below=of start] (solve) {
    \textbf{计算高斯点处的通量}
  };

  \node[mybox, below=0.9cm of solve] (robust1) {
    \textbf{使用保正限制器提升算法鲁棒性}
  };

  \node[mybox, below=0.9cm of robust1] (robust2) {
    \textbf{使用近似黎曼求解器计算数值通量}
  };

  \node[mybox, below=0.9cm of robust2] (update) {
    \textbf{更新流场和当前计算时间}\\[2pt]
    $t^{n+1} = t^n + \Delta t$，\quad
    $u^{n+1} = u^n + \Delta u$
  };

  \node[mydiamond, below=0.9cm of update] (check) {
    是否达到终止时间 $T$?
  };

  \node[mydiamond, below=0.9cm of check] (stop) {停止计算};

  % ── 右侧说明框 ───────────────────────────────────
  \node[sidegreen, text width=4.2cm, right=1.6cm of solve] (note1) {
    \textbf{求解步骤}\\[3pt]
    1. 选择合适的函数空间\\[3pt]
    2. 优化形状参数\\[3pt]
    3. WENO 插值
  };

  \node[sidegreen, text width=4.2cm, right=1.6cm of robust2] (note2) {
    \textbf{近似黎曼求解器选择}\\[3pt]
    $\bullet$ LLF Flux（简便）\\[3pt]
    $\bullet$ HLLC Flux（耗散小）\upcite{ZHANG2025116467}
  };

  % ── 左侧循环初值框 ───────────────────────────────
  \node[sideleft, left=2.2cm of robust2] (reinit) {
    将 $t^{n+1}$ 时刻的\\数值解\\作为初值解
  };

  % ── 主流程箭头 ───────────────────────────────────
  \draw[myarrow] (init)    -- (start);
  \draw[myarrow] (start)   -- (solve);
  \draw[myarrow] (solve)   -- (robust1);
  \draw[myarrow] (robust1) -- (robust2);
  \draw[myarrow] (robust2) -- (update);
  \draw[myarrow] (update)  -- (check);
  \draw[myarrow] (check)   -- node[right, font=\normalsize, red!70!black] {是} (stop);

  % 否 → 左侧回路
  \draw[myarrow] (check.west)
    -| node[left, font=\normalsize, red!70!black] {否}
    (reinit.south);
  \draw[myarrow] (reinit.north)
    |- (start.west);

  % ── 右侧绿色箭头 ─────────────────────────────────
  \draw[greenarrow] (note1.west) -- (solve.east);
  \draw[greenarrow] (note2.west) -- (robust2.east);

\end{tikzpicture}
}